\documentclass[11pt]{article}
\usepackage[a4paper,margin=2cm]{geometry}
\usepackage{amsmath,amssymb}
\usepackage{xcolor}
\usepackage{fancyhdr}
\usepackage{enumitem}

\definecolor{copper}{HTML}{8C4A2F}
\definecolor{iron}{HTML}{262B33}
\definecolor{muted}{HTML}{6B6B6B}

\pagestyle{fancy}
\fancyhf{}
\lhead{\footnotesize\color{muted} SP-03 \textbar{} Transformer III: Voltage Regulation \textbar{} EM-MTE notes}
\rhead{\footnotesize\color{muted} ELC2104 \textbar{} MTE}
\cfoot{\footnotesize\color{muted} \thepage}

\setlength{\parindent}{0pt}
\setlength{\parskip}{5pt}
\renewcommand{\familydefault}{\sfdefault}

\newcommand{\qhead}[2]{\subsection*{\color{copper}#1\quad\color{iron}#2}}
\newcommand{\tagline}[1]{\textbf{\color{copper}#1}}
\newcommand{\keydist}[1]{\vspace{2pt}\textbf{\color{red}KEY DISTINCTION.} #1}
\newcommand{\drawline}[1]{\textbf{\color{muted}[DRAW]} #1}
\newcommand{\srcwrong}[1]{\textbf{\color{red}[SRC WRONG]} #1}
\newcommand{\trap}[1]{\vspace{2pt}\textbf{\color{red}[TRAP]} #1}

\begin{document}

\begin{center}
{\footnotesize\color{copper} ELC2104 \textbar{} ELECTRICAL MACHINES \textbar{} MTE}\\[6pt]
{\LARGE\bfseries\color{iron} SP-03: Transformer III}\\[2pt]
{\Large\color{copper} Voltage Regulation: Exact, Approximate, and the Percent-Drop Shortcut}\\[6pt]
{\footnotesize\color{muted} Question-driven notes, dense-point standard. Covers T.6 of the locked syllabus.}
\end{center}

\vspace{2pt}
\tagline{QUESTION SHAPES IN THIS FILE:} Q6 (explain regulation + derive at different power factors), Q8 (derive the approximate expression, lagging and leading). Numerics: WE-08, WE-09, D1-01, D1-02 (regulation and percent-drop families).
\vspace{2pt}

\keydist{Regulation is a statement about the \emph{internal drops}, not the load: same primary voltage, secondary open versus secondary loaded. Lagging loads make the output sag. Leading loads can make the output \emph{rise} past its no-load value (negative regulation), and there is one specific leading power factor where regulation is exactly zero.}

\hrulefill

% ============================================================ Q6
\qhead{Q6 (verbatim):}{Explain voltage regulation of a transformer and derive the expression for voltage regulation at different power factors.}

\tagline{WHAT TO WRITE:} definition (with the test condition: constant primary voltage) \textrightarrow{} why the drop exists \textrightarrow{} the phasor picture in words \textrightarrow{} the exact derivation \textrightarrow{} the percentage form and the denominator convention \textrightarrow{} the three power-factor cases.

\tagline{DEFINITION AND PHYSICS (points 1 to 4).}
\begin{enumerate}[label=\arabic*.]
\item \textbf{Voltage regulation} is the change in secondary terminal voltage when the load is thrown off, with the \emph{primary supply voltage held constant}. Practically: the voltmeter at the secondary reads lower (usually) at load than at no load, because of the internal voltage drops.
\item The cause is the two series effects in the equivalent circuit: the resistive drop ($I_2 R_{\text{eq}}$, in phase with the current) and the leakage-reactive drop ($I_2 X_{\text{eq}}$, leading the current by 90$^\circ$). Together they make the internal induced emf $E_2$ larger in magnitude than the terminal voltage $V_2$.
\item Definition in symbols (absolute regulation):
\[
\text{Regulation} = E_2 - V_2
\]
where $V_2$ = secondary terminal voltage \emph{on load} and $E_2$ = secondary terminal voltage at \emph{no load} with the same primary voltage (equal to the induced emf, since at no load there is no drop).
\item In percentage form (the exam's usual ask):
\[
\%\,\text{VR} = \frac{E_2 - V_2}{E_2} \times 100
\qquad\text{(the deck's convention: divide by the no-load voltage)}
\]
\trap{the denominator war (H.3): the handwritten notes use $\%\,\text{VR} = \dfrac{E_2 - V_2}{V_2} \times 100$ (divide by the load voltage, treated as the rated secondary voltage), and some textbooks divide by the rated secondary voltage too. The numeric difference is small (a few percent of a few percent) but the marker wants \emph{one} convention used consistently and \emph{named}. Exam rule for this course: follow the deck convention (divide by the no-load value $E_2 = K V_1$), and state ``regulation referred to the no-load secondary voltage'' in one line. If you compute with the other one anyway, say so.}
\end{enumerate}

\tagline{DERIVATION OF THE EXACT EXPRESSION (lagging load).}

\textbf{Step 1: set the phasor frame.} Take $V_2$ along the reference axis. The load current lags $V_2$ by the load angle $\phi$, so its components are $I_2\cos\phi$ (along $V_2$) and $I_2\sin\phi$ (behind $V_2$).

\textbf{Step 2: place the drops.} The copper drop $I_2 R_{\text{eq}}$ is in phase with $I_2$; the reactive drop $I_2 X_{\text{eq}}$ leads $I_2$ by 90$^\circ$. Resolve both onto the $V_2$ axis and the perpendicular axis:
\begin{align*}
\text{along } V_2:&\quad V_2\cos\phi + I_2 R_{\text{eq}} \\
\text{perpendicular, in the direction of } E_2:&\quad V_2\sin\phi + I_2 X_{\text{eq}}
\end{align*}
(the reactive drop, leading the current, lands in the same quadrant as the $V_2\sin\phi$ component for a lagging load)

\textbf{Step 3: Pythagoras gives the exact result:}
\[
E_2 = \sqrt{\left(V_2\cos\phi + I_2 R_{\text{eq}}\right)^{2} + \left(V_2\sin\phi + I_2 X_{\text{eq}}\right)^{2}}
\qquad\text{(lagging power factor)}
\]

\textbf{Step 4: leading power factor.} The current leads $V_2$ by $\phi$, so the reactive drop subtracts from the $V_2\sin\phi$ component:
\[
E_2 = \sqrt{\left(V_2\cos\phi + I_2 R_{\text{eq}}\right)^{2} + \left(V_2\sin\phi - I_2 X_{\text{eq}}\right)^{2}}
\qquad\text{(leading power factor)}
\]

\textbf{Step 5: unity power factor.} $\sin\phi = 0$, so only the resistive term survives:
\[
E_2 = \sqrt{\left(V_2 + I_2 R_{\text{eq}}\right)^{2} + \left(I_2 X_{\text{eq}}\right)^{2}}
\]

\textbf{Step 6: regulation at each case} is then $\%\,\text{VR} = \dfrac{E_2 - V_2}{E_2} \times 100$ with $E_2$ from the matching line above.

\tagline{THE THREE CASES, SUMMARISED (write this table in words).}
\begin{enumerate}[label=\arabic*.]
\item \textbf{Lagging pf:} regulation is positive and \emph{largest}. Both drop components push $E_2$ above $V_2$. Worse power factor (smaller $\cos\phi$) means worse regulation.
\item \textbf{Unity pf:} regulation is positive but \emph{smaller}: only $I_2 R_{\text{eq}}$ and the quadrature bookkeeping of $I_2 X_{\text{eq}}$ enter the exact formula.
\item \textbf{Leading pf:} regulation falls, can pass through \emph{zero}, and can go \emph{negative} (the load voltage rises above the no-load value: the capacitive current's magnetising effect supports the flux). Deep leading power factor is the classic ``voltage rises at the receiving end'' case.
\end{enumerate}

\drawline{phasor diagram for the lagging case: $V_2$ reference, $I_2$ behind it by $\phi$, $I_2 R_{\text{eq}}$ drawn along $I_2$, $I_2 X_{\text{eq}}$ perpendicular to $I_2$ (leading it), $E_2$ closing the polygon. Repeat the same figure with the current leading (and the reactive drop subtracting) for the leading case. Label checklist: $V_2$, $E_2$, $I_2$, $\phi$, both drop segments, the right angles.}

\keydist{The definition's condition is \emph{constant primary voltage}. If the question says the primary sags too, regulation no longer describes the result: the source impedance joins the story. Keep the definition clean and the derivation clean.}

% ============================================================ Q8
\qhead{Q8 (verbatim):}{Derive the approximate expression for voltage regulation of a transformer for lagging and leading power factors.}

\tagline{WHAT TO WRITE:} start from the exact squared relation \textrightarrow{} expand \textrightarrow{} argue which terms are negligible \textrightarrow{} the approximate drop \textrightarrow{} the percentage form with $\%R$ and $\%X$ \textrightarrow{} the shortcut corollaries.

\tagline{DERIVATION.}

\textbf{Step 1: the exact relation, squared (lagging).} From Q6:
\[
E_2^{2} = \left(V_2\cos\phi + I_2 R_{\text{eq}}\right)^{2} + \left(V_2\sin\phi + I_2 X_{\text{eq}}\right)^{2}
\]

\textbf{Step 2: expand.}
\[
E_2^{2} = V_2^{2}\cos^{2}\phi + 2 V_2 I_2 R_{\text{eq}}\cos\phi + I_2^{2} R_{\text{eq}}^{2} + V_2^{2}\sin^{2}\phi + 2 V_2 I_2 X_{\text{eq}}\sin\phi + I_2^{2} X_{\text{eq}}^{2}
\]
\[
E_2^{2} = V_2^{2} + 2 V_2 I_2\left(R_{\text{eq}}\cos\phi + X_{\text{eq}}\sin\phi\right) + I_2^{2}\left(R_{\text{eq}}^{2} + X_{\text{eq}}^{2}\right)
\]
using $\cos^{2}\phi + \sin^{2}\phi = 1$ and $R_{\text{eq}}^{2} + X_{\text{eq}}^{2} = Z_{\text{eq}}^{2}$.

\textbf{Step 3: the approximation argument (this sentence earns marks).} In a normal transformer the short-circuit impedance is 5 to 10 percent, so the drop terms $I_2 Z_{\text{eq}}$ are small against $V_2$. The squared drop $I_2^{2} Z_{\text{eq}}^{2}$ is second-order small compared with the product term $2 V_2 I_2 Z_{\text{eq}}$ and is neglected (equivalently: use $\sqrt{1 + \epsilon} \approx 1 + \dfrac{\epsilon}{2}$ with $\epsilon \ll 1$).

\textbf{Step 4: the approximate drop.} Then $E_2 \approx V_2 + I_2 R_{\text{eq}}\cos\phi + I_2 X_{\text{eq}}\sin\phi$, so:
\[
\boxed{\;E_2 - V_2 \approx I_2 R_{\text{eq}}\cos\phi + I_2 X_{\text{eq}}\sin\phi\;}
\qquad\text{(lagging)}
\]
For a \textbf{leading} load the reactive term subtracts:
\[
\boxed{\;E_2 - V_2 \approx I_2 R_{\text{eq}}\cos\phi - I_2 X_{\text{eq}}\sin\phi\;}
\qquad\text{(leading)}
\]

\textbf{Step 5: percentage form.} Divide by the chosen reference voltage and multiply by 100. With the notes convention (divide by $V_2$):
\[
\%\,\text{VR} = \frac{I_2 R_{\text{eq}}}{V_2}\cos\phi \times 100 \;\pm\; \frac{I_2 X_{\text{eq}}}{V_2}\sin\phi \times 100
\]
Defining the \textbf{percent drops at full load}:
\[
\%\,R_{\text{drop}} = \frac{I_{2,\text{FL}} R_{\text{eq}}}{V_2} \times 100
\qquad\qquad
\%\,X_{\text{drop}} = \frac{I_{2,\text{FL}} X_{\text{eq}}}{V_2} \times 100
\]
the famous compact line (the exam's ``express regulation in terms of percent drops'' answer):
\[
\boxed{\;\%\,\text{VR} = \left(\%\,R_{\text{drop}}\right)\cos\phi \;\pm\; \left(\%\,X_{\text{drop}}\right)\sin\phi\;}
\]
at fractional load $x$ every drop scales linearly with $x$ (this is load current, not loss: the loss scales as $x^{2}$, the drop as $x$).

\tagline{ZERO REGULATION (derive when asked, quote always).} Set the leading expression to zero:
\[
R_{\text{eq}}\cos\phi = X_{\text{eq}}\sin\phi
\qquad\Longrightarrow\qquad
\tan\phi = \frac{R_{\text{eq}}}{X_{\text{eq}}}
\qquad\text{(leading, i.e. capacitive load)}
\]
At exactly this leading power factor the load voltage equals the no-load voltage. For a more capacitive load than this, regulation turns negative.

\tagline{SHORTCUT COROLLARIES (the percent-impedance family).}
\begin{enumerate}[label=\arabic*.]
\item The percent resistive drop equals the copper loss as a percentage of the kVA rating:
\[
\%\,R_{\text{drop}} = \frac{P_{\text{cu,FL}}}{S_{\text{rated}}} \times 100
\]
(because $I_2 R_{\text{eq}} = \dfrac{I_2^{2} R_{\text{eq}}}{I_2} = \dfrac{P_{\text{cu,FL}}}{I_2}$ and $S = V_2 I_2$: dividing gives the identity).
\item The percent impedance drop comes straight from the SC test voltage at rated current:
\[
\%\,Z_{\text{drop}} = \frac{V_{\text{sc}}}{V_{\text{rated, supplied side}}} \times 100
\qquad\Longrightarrow\qquad
\%\,X_{\text{drop}} = \sqrt{\left(\%\,Z_{\text{drop}}\right)^{2} - \left(\%\,R_{\text{drop}}\right)^{2}}
\]
\item With only percent numbers given (the D1-02 and WE-09 shape), no circuit is needed at all: plug straight into the compact line.
\end{enumerate}

\trap{do not confuse the drop scaling with the loss scaling: drops scale linearly with the load fraction $x$, copper losses scale as $x^{2}$. One is regulation, the other is efficiency (SP-04). Mixing them is a two-mistake freebie.}

\drawline{the impedance triangle in percent form (horizontal $\%\,R$, vertical $\%\,X$, hypotenuse $\%\,Z$) and the zero-regulation phasor (leading current sized so that $I_2 X_{\text{eq}}$ exactly cancels the $V_2\sin\phi$ build-up).}

\keydist{The exact expression is for computing. The approximate expression is for \emph{explaining} (it shows the two drop mechanisms add or subtract by power factor). The percent form is for fast exams. Know all three and say which one the numbers came from.}

% ============================================================ NUMERICS
\hrulefill
\subsection*{\color{copper}WORKED NUMERICS: REGULATION FAMILIES (WE-08, WE-09, D1-01, D1-02)}

\tagline{WE-08 (full method).} A 40 kVA, 6600/250 V transformer has $R_1 = 10\ \Omega$ (primary), $R_2 = 0.02\ \Omega$ (secondary), and total equivalent leakage reactance referred to the primary $X_{\text{eq,1}} = 35\ \Omega$. Find the full-load voltage regulation at 0.8 power factor lagging.

\textbf{Givens:} $S = 40$ kVA, $V_1 = 6\,600$ V, $V_2 = 250$ V, $R_1 = 10\ \Omega$, $R_2 = 0.02\ \Omega$, $X_{\text{eq,1}} = 35\ \Omega$, $\cos\phi = 0.8$ lagging ($\sin\phi = 0.6$).

\textbf{Solution. Work on the secondary side (it is where regulation lives).}
\[
a = \frac{6\,600}{250} = 26.4
\qquad\Longrightarrow\qquad
R_1' = \frac{R_1}{a^{2}} = 10 \times \left(\frac{250}{6\,600}\right)^{2} = 0.01435\ \Omega
\qquad
X_{\text{eq,2}} = 35 \times \left(\frac{250}{6\,600}\right)^{2} = 0.05022\ \Omega
\]
\[
R_{\text{eq,2}} = 0.01435 + 0.02 = 0.03435\ \Omega
\qquad\qquad
I_{2,\text{FL}} = \frac{40\,000}{250} = 160\ \text{A}
\]
\[
E_2 - V_2 \approx I_2\left(R_{\text{eq,2}}\cos\phi + X_{\text{eq,2}}\sin\phi\right) = 160\,(0.02748 + 0.03013) = 9.22\ \text{V}
\]
\[
\%\,\text{VR} = \frac{9.22}{250} \times 100 = 3.69\%
\qquad\text{(notes convention: divide by the load voltage)}
\]
One line for the deck convention (divide by the no-load voltage $E_2 = 250 + 9.22 = 259.22$ V): $\%\,\text{VR} = \dfrac{9.22}{259.22} \times 100 = 3.56\%$. \textbf{Answer:} regulation drop = 9.22 V, $\%\,\text{VR} = 3.69\%$ (or 3.56\% on the no-load denominator: whichever you use, name it).

\vspace{2pt}

\tagline{WE-09 (percent-drop method, errata E9).} The ohmic drop of a transformer is 1 percent of the rated voltage and the reactance drop is 5 percent. Find the regulation at (1) 0.8 power factor lagging, (2) 0.8 power factor leading.

\textbf{Solution (no circuit needed, straight into the compact line).}
\[
\%\,\text{VR} = \left(\%\,R_{\text{drop}}\right)\cos\phi \pm \left(\%\,X_{\text{drop}}\right)\sin\phi
\qquad\text{with}\quad \cos\phi = 0.8,\;\; \sin\phi = 0.6
\]
\[
\text{(1) lagging:}\quad \%\,\text{VR} = 1 \times 0.8 + 5 \times 0.6 = 0.8 + 3 = 3.8\%
\]
\[
\text{(2) leading:}\quad \%\,\text{VR} = 1 \times 0.8 - 5 \times 0.6 = 0.8 - 3 = -2.2\%
\]
\textbf{Answer:} 3.8\% lagging, $-2.2\%$ leading. The negative sign means the load voltage \emph{rises} about 2.2 percent above the no-load value. \srcwrong{the source labels these quantities as ``volts''. They are percentages (errata E9): the drops were given in percent, so the results are in percent.}

\vspace{2pt}

\tagline{D1-01 (percent drops at three power factors + the zero-regulation angle).} The ohmic drop of a transformer is 1 percent of output and the reactance drop is 5 percent of the voltage. Find the regulation at (a) 0.8 lagging, (b) unity, (c) 0.8 leading. Then find the leading power factor at which regulation is exactly zero.

\textbf{Solution.}
\[
\text{(a)}\;\; 1 \times 0.8 + 5 \times 0.6 = 3.8\%
\qquad
\text{(b)}\;\; 1 \times 1 + 5 \times 0 = 1\%
\qquad
\text{(c)}\;\; 1 \times 0.8 - 5 \times 0.6 = -2.2\%
\]
\textbf{Zero regulation:}
\[
\tan\phi = \frac{\%\,R_{\text{drop}}}{\%\,X_{\text{drop}}} = \frac{1}{5}
\qquad\Longrightarrow\qquad
\phi = 11.31^\circ \text{ leading}
\qquad\Longrightarrow\qquad
\cos\phi = 0.981 \text{ leading}
\]
Check: $1 \times 0.9806 - 5 \times 0.1961 = 0.9806 - 0.9806 = 0$ $\checkmark$.
\textbf{Answer:} (a) 3.8\%, (b) 1\%, (c) $-2.2\%$; zero regulation at 0.981 leading.

\vspace{2pt}

\tagline{D1-02 (the full circuit variant).} A 10 kVA, 2000/400 V transformer has $R_1 = 5.5\ \Omega$, $X_1 = 12\ \Omega$, $R_2 = 0.2\ \Omega$, $X_2 = 0.45\ \Omega$. Find the approximate secondary voltage at full load, 0.8 power factor lagging, when the primary is at 2000 V.

\textbf{Solution.}
\[
a = \frac{2\,000}{400} = 5
\qquad\Longrightarrow\qquad
R_{\text{eq,1}} = 5.5 + 25 \times 0.2 = 10.5\ \Omega
\qquad
X_{\text{eq,1}} = 12 + 25 \times 0.45 = 23.25\ \Omega
\]
Refer to the secondary (divide by $a^{2} = 25$) and take the full-load secondary current:
\[
R_{\text{eq,2}} = \frac{10.5}{25} = 0.42\ \Omega
\qquad
X_{\text{eq,2}} = \frac{23.25}{25} = 0.93\ \Omega
\qquad
I_{2,\text{FL}} = \frac{10\,000}{400} = 25\ \text{A}
\]
The no-load secondary voltage with 2000 V applied is $\dfrac{2\,000}{5} = 400$ V, so:
\[
E_2 - V_2 \approx 25\,(0.42 \times 0.8 + 0.93 \times 0.6) = 25\,(0.336 + 0.558) = 22.35\ \text{V}
\]
\[
V_2 \approx 400 - 22.35 = 377.65\ \text{V} \approx 377.7\ \text{V}
\qquad\qquad
\%\,\text{VR} = \frac{22.35}{400} \times 100 = 5.59\%
\]
\textbf{Answer:} about 377.7 V at full load 0.8 lagging (regulation 5.59\% on the no-load denominator here).

% ============================================================ SHEET
\hrulefill
\subsection*{\color{copper}FORMULA SHEET FOR THIS FILE}
\begin{align*}
\text{definition:}\quad \text{Regulation} &= E_2 - V_2
& \%\,\text{VR} &= \frac{E_2 - V_2}{E_2} \times 100 \;\; (\text{deck}) \\[4pt]
\text{exact, lagging:}\quad E_2 &= \sqrt{(V_2\cos\phi + I_2 R_{\text{eq}})^{2} + (V_2\sin\phi + I_2 X_{\text{eq}})^{2}} \\[4pt]
\text{exact, leading:}\quad E_2 &= \sqrt{(V_2\cos\phi + I_2 R_{\text{eq}})^{2} + (V_2\sin\phi - I_2 X_{\text{eq}})^{2}} \\[4pt]
\text{approx:}\quad E_2 - V_2 &\approx I_2 R_{\text{eq}}\cos\phi \pm I_2 X_{\text{eq}}\sin\phi \\[4pt]
\%\,\text{VR} &= (\%\,R_{\text{drop}})\cos\phi \pm (\%\,X_{\text{drop}})\sin\phi \\[4pt]
\%\,R_{\text{drop}} &= \frac{P_{\text{cu,FL}}}{S} \times 100
& \%\,X_{\text{drop}} &= \sqrt{(\%\,Z_{\text{drop}})^{2} - (\%\,R_{\text{drop}})^{2}} \\[4pt]
\text{zero regulation:}\quad \tan\phi &= \frac{R_{\text{eq}}}{X_{\text{eq}}}
& &\text{(leading pf)}
\end{align*}

\subsection*{\color{copper}CLOSED-BOOK CHECK (cover the file, answer out loud)}
\begin{enumerate}[label=\arabic*.]
\item Define regulation with its test condition and both percentage conventions (and say which one you will use and why).
\item Derive the exact lagging expression: the phasor resolution in words, then the two squared terms.
\item Get the approximate expression from the exact one and state the size argument that licenses the neglect.
\item State the three cases (lag / unity / lead) and what each does to the regulation, including the sign.
\item Derive the zero-regulation power factor and say what happens on either side of it.
\item From memory: WE-08's drop and both percentage conventions; D1-02's $V_2$.
\end{enumerate}

\subsection*{\color{copper}MEMORY DUMP (write cold on blank paper)}
Regulation = $E_2 - V_2$ at constant $V_1$. Exact lagging: $E_2 = \sqrt{(V_2\cos\phi + I_2 R_{\text{eq}})^{2} + (V_2\sin\phi + I_2 X_{\text{eq}})^{2}}$; leading flips the $X$ sign. Approximate: $E_2 - V_2 \approx I_2 R_{\text{eq}}\cos\phi \pm I_2 X_{\text{eq}}\sin\phi$ (neglect $(I_2 Z_{\text{eq}})^{2}$). Percent line: $\%\,\text{VR} = (\%\,R)\cos\phi \pm (\%\,X)\sin\phi$; $\%\,R = \dfrac{P_{\text{cu,FL}}}{S} \times 100$; $\%\,X = \sqrt{\%\,Z^{2} - \%\,R^{2}}$. Lag: positive and worst. Unity: small positive. Lead: shrinks, zero at $\tan\phi = \dfrac{R_{\text{eq}}}{X_{\text{eq}}}$, then negative (voltage rises). Drops scale as $x$, copper loss as $x^{2}$. Name the denominator convention every time.

\end{document}
